\documentclass[11pt]{article}
\usepackage{amsmath,amssymb,booktabs,geometry,array}
\geometry{left=2.5cm,right=2.5cm,top=2.5cm,bottom=2.5cm}
\setlength{\parskip}{6pt}\setlength{\parindent}{0pt}
\title{Correlation structure of the \texttt{covar} and \texttt{orig}
data-generating processes}
\author{pmsimstats team}
\date{\today}

\newcommand{\bm}{\mathrm{BM}}
\newcommand{\bl}{\mathrm{BL}}

\begin{document}
\maketitle

\section{Variable ordering and block partition}

Both implementations draw each participant's trajectory from a single
multivariate normal over the same ordered variable list. With $n$
measurement occasions at cumulative weeks
$w_1 < w_2 < \cdots < w_n$, the vector is

\[
X = \bigl(\,\bm,\; \bl,\;
TV_1,\dots,TV_n,\;
PB_1,\dots,PB_n,\;
BR_1,\dots,BR_n \,\bigr)^{\!\top},
\qquad \dim X = 2 + 3n .
\]

The ordering is factor-major: all $n$ occasions of one factor are
contiguous, rather than all three factors at one occasion. Writing
$A_c$ for the $n \times n$ within-factor block of factor
$c \in \{tv, pb, br\}$ and $C_{cc'}$ for the $n \times n$
cross-factor block, the correlation matrix partitions as

\[
R \;=\;
\begin{pmatrix}
1 & 0 & 0^{\top} & 0^{\top} & b^{\top} \\[2pt]
0 & 1 & 0^{\top} & 0^{\top} & 0^{\top} \\[2pt]
0 & 0 & A_{tv}   & C_{tv,pb} & C_{tv,br} \\[2pt]
0 & 0 & C_{pb,tv} & A_{pb}   & C_{pb,br} \\[2pt]
b & 0 & C_{br,tv} & C_{br,pb} & A_{br}
\end{pmatrix}.
\]

Two features hold in both implementations. The baseline term $\bl$ is
uncorrelated with every other variable, so its row and column are
those of the identity. The biomarker $\bm$ couples only to the
response factor $BR$, through the vector $b \in \mathbb{R}^{n}$; its
correlations with $TV$ and $PB$ are zero. The two implementations
differ in the interior of $A_c$, in the interior of $C_{cc'}$, and in
$b$.

\section{Within-factor blocks $A_c$}

Let $\rho_c$ denote the autocorrelation parameter of factor $c$
(\texttt{c.tv}, \texttt{c.pb}, \texttt{c.br} in the code).

\paragraph{\texttt{orig}: compound symmetry.}
A single constant is assigned to every pair of occasions, irrespective
of their separation:
\[
[A_c]_{ij} =
\begin{cases}
1, & i = j\\
\rho_c, & i \neq j
\end{cases}
\qquad\Longleftrightarrow\qquad
A_c = (1-\rho_c) I_n + \rho_c J_n ,
\]
with $J_n$ the $n \times n$ matrix of ones. For $n = 4$,
\[
A_c^{\text{orig}} =
\begin{pmatrix}
1 & \rho_c & \rho_c & \rho_c\\
\rho_c & 1 & \rho_c & \rho_c\\
\rho_c & \rho_c & 1 & \rho_c\\
\rho_c & \rho_c & \rho_c & 1
\end{pmatrix}.
\]
The block has two distinct eigenvalues, $1 + (n-1)\rho_c$ once and
$1 - \rho_c$ with multiplicity $n-1$, so positive definiteness
requires $-1/(n-1) < \rho_c < 1$. The upper bound on admissible
correlation therefore tightens as occasions are added.

\paragraph{\texttt{covar}: AR(1) in calendar time.}
The correlation decays with the elapsed time between occasions, using
the cumulative week gap rather than the index difference:
\[
[A_c]_{ij} = \rho_c^{\,|w_i - w_j|}.
\]
For $n = 4$, writing $d_{ij} = |w_i - w_j|$,
\[
A_c^{\text{covar}} =
\begin{pmatrix}
1 & \rho_c^{d_{12}} & \rho_c^{d_{13}} & \rho_c^{d_{14}}\\
\rho_c^{d_{12}} & 1 & \rho_c^{d_{23}} & \rho_c^{d_{24}}\\
\rho_c^{d_{13}} & \rho_c^{d_{23}} & 1 & \rho_c^{d_{34}}\\
\rho_c^{d_{14}} & \rho_c^{d_{24}} & \rho_c^{d_{34}} & 1
\end{pmatrix}.
\]
Because the exponent uses calendar time, unequally spaced designs
produce unequal decay across consecutive occasions. In the Hybrid
design, whose occasions fall at weeks
$4, 8, 9, 10, 11, 12, 16, 20$, consecutive gaps of one week during
the discontinuation phase yield much higher adjacent correlations
than the four-week gaps at either end.

\section{Cross-factor blocks $C_{cc'}$}

Let $c_1$ denote the same-occasion cross-factor correlation
(\texttt{c.cf1t}) and $c_{\times}$ the different-occasion cross-factor
correlation (\texttt{c.cfct}).

\paragraph{\texttt{orig}.} Both are constants:
\[
[C_{cc'}]_{ij} =
\begin{cases}
c_1, & i = j\\
c_{\times}, & i \neq j
\end{cases}
\qquad\Longleftrightarrow\qquad
C_{cc'} = c_{\times} J_n + (c_1 - c_{\times}) I_n .
\]

\paragraph{\texttt{covar}.} The same-occasion entry is unchanged, but
the different-occasion entries inherit the AR(1) decay:
\[
[C_{cc'}]_{ij} =
\begin{cases}
c_1, & i = j\\
c_{\times}\,\rho^{\,|w_i - w_j|}, & i \neq j .
\end{cases}
\]

\section{Biomarker coupling vector $b$}

The vector $b$ carries the biomarker-treatment interaction. This is
the only place in $R$ where the two implementations differ in
mechanism rather than in decay profile.

\paragraph{\texttt{covar}.} The entry is gated on the drug indicator
and decays with time since discontinuation $t_{sd,p}$ at rate
$\lambda = \ln 2 / t_{1/2}$:
\[
b_p =
\begin{cases}
c_{bm}, & \text{occasion } p \text{ on drug}\\
c_{bm}\,e^{-\lambda t_{sd,p}}, & \text{occasion } p \text{ off drug} .
\end{cases}
\]

\paragraph{\texttt{orig}.} The entry is gated on whether the response
mean $\mu_{BR,p}$ is nonzero, and takes only two values:
\[
b_p =
\begin{cases}
c_{bm}, & \mu_{BR,p} \neq 0\\
0, & \mu_{BR,p} = 0 .
\end{cases}
\]

Because carryover keeps $\mu_{BR,p}$ nonzero after discontinuation,
the position of this step depends on the carryover half-life. At
$t_{1/2} = 0$ the step coincides with the drug indicator and the two
vectors are equal entry for entry. As $t_{1/2}$ grows they separate in
opposite directions: the \texttt{covar} entries decay toward zero
while the \texttt{orig} entries rise to $c_{bm}$. At
$t_{1/2} = 1.0$ in the Hybrid design, $b^{\text{orig}} = c_{bm}
\mathbf{1}_n$, so the vector carries no on-drug versus off-drug
contrast at all.

\section{A worked example}

To make the three components concrete, take four occasions drawn from
the Hybrid design at cumulative weeks
$w = (10, 11, 12, 16)$, with the first and fourth on drug and the
second and third off it, so that
$t_{sd} = (0, 1, 2, 0)$. Set $\rho_{br} = 0.6$, $\rho_{pb} = 0.4$,
$c_1 = 0.3$, $c_{\times} = 0.1$, $c_{bm} = 0.45$, and a carryover
half-life of one week, giving $\lambda = \ln 2 \approx 0.693$. The
matrix of week gaps is
\[
\bigl(|w_i - w_j|\bigr) =
\begin{pmatrix}
0 & 1 & 2 & 6\\
1 & 0 & 1 & 5\\
2 & 1 & 0 & 4\\
6 & 5 & 4 & 0
\end{pmatrix}.
\]
The unequal spacing matters: occasions one to three are a week apart,
while the fourth sits four weeks after the third.

\subsection{Within-factor block $A_{br}$}

\[
A_{br}^{\text{orig}} =
\begin{pmatrix}
1.0000 & 0.6000 & 0.6000 & 0.6000 \\
0.6000 & 1.0000 & 0.6000 & 0.6000 \\
0.6000 & 0.6000 & 1.0000 & 0.6000 \\
0.6000 & 0.6000 & 0.6000 & 1.0000
\end{pmatrix},
\qquad
A_{br}^{\text{covar}} =
\begin{pmatrix}
1.0000 & 0.6000 & 0.3600 & 0.0467 \\
0.6000 & 1.0000 & 0.6000 & 0.0778 \\
0.3600 & 0.6000 & 1.0000 & 0.1296 \\
0.0467 & 0.0778 & 0.1296 & 1.0000
\end{pmatrix}.
\]

Under \texttt{orig} every pair of occasions carries $0.6$, so the
week-apart pair $(1,2)$ and the six-weeks-apart pair $(1,4)$ are
treated as equally associated. Under \texttt{covar} the same two pairs
carry $0.6$ and $0.0467$, a factor of thirteen apart. The consequence
for the fourth occasion is the striking one: it is nearly independent
of the first three under AR(1), and as strongly tied to them as they
are to each other under compound symmetry.

\subsection{Cross-factor block $C_{tv,pb}$}

\[
C_{tv,pb}^{\text{orig}} =
\begin{pmatrix}
0.3000 & 0.1000 & 0.1000 & 0.1000 \\
0.1000 & 0.3000 & 0.1000 & 0.1000 \\
0.1000 & 0.1000 & 0.3000 & 0.1000 \\
0.1000 & 0.1000 & 0.1000 & 0.3000
\end{pmatrix},
\qquad
C_{tv,pb}^{\text{covar}} =
\begin{pmatrix}
0.3000 & 0.0400 & 0.0160 & 0.0004 \\
0.0400 & 0.3000 & 0.0400 & 0.0010 \\
0.0160 & 0.0400 & 0.3000 & 0.0026 \\
0.0004 & 0.0010 & 0.0026 & 0.3000
\end{pmatrix}.
\]

The diagonal, which carries the same-occasion cross-factor
correlation $c_1$, is identical in both. The off-diagonal entries
decay under \texttt{covar} and do not under \texttt{orig}, and by the
fourth occasion the difference is two orders of magnitude.

A caution attaches to the \texttt{covar} block. In the implementation
the decay rate applied to a cross-factor pair is taken from whichever
factor the outer loop is currently visiting, and because each
unordered pair is visited in both orders, the later write survives.
The rate appearing in $C_{tv,pb}$ is therefore determined by loop
order rather than by any modeling decision. The values above use
$\rho_{pb}$ on that basis. This is an implementation artifact and not
a property of the AR(1) specification.

\subsection{Biomarker vector $b$}

At a carryover half-life of one week,
\[
b^{\text{covar}} = (0.4500,\, 0.2250,\, 0.1125,\, 0.4500)^{\top},
\qquad
b^{\text{orig}} = (0.4500,\, 0.4500,\, 0.4500,\, 0.4500)^{\top}.
\]
The two off-drug entries halve and quarter under \texttt{covar},
tracking $e^{-\lambda t_{sd}}$, whereas under \texttt{orig} they take
the full $c_{bm}$, because carryover has left the response mean
nonzero and the step therefore never fires.

Removing carryover collapses the difference entirely:
\[
b^{\text{covar}} = b^{\text{orig}} =
(0.4500,\, 0.0000,\, 0.0000,\, 0.4500)^{\top}
\qquad (t_{1/2} = 0).
\]

This pair of displays is the essential comparison. The vector $b$ is
the only channel through which the biomarker-treatment interaction
enters the covariance, and an interaction is identified by $b$ varying
between on-drug and off-drug occasions. At $t_{1/2} = 0$ the two
implementations produce the same vector. At any positive half-life
$b^{\text{orig}}$ is constant, so it carries no on-drug versus
off-drug contrast, while $b^{\text{covar}}$ retains one. Across the
grid $t_{1/2} \in \{0, 0.5, 1.0\}$ used in this paper, \texttt{orig}
therefore supplies one cell identical to \texttt{covar} and two in
which its interaction channel is degenerate.

\section{The three data-generating processes}

Organized by where each process writes the biomarker-treatment
interaction into the multivariate normal $(\mu, \Sigma)$. All three
share the same variable ordering and block partition of Section
1, and differ only in the entries listed.

\begin{center}
\small
\begin{tabular}{@{}
  >{\raggedright\arraybackslash}p{3.5cm}
  >{\raggedright\arraybackslash}p{2.5cm}
  >{\raggedright\arraybackslash}p{3.5cm}
  >{\raggedright\arraybackslash}p{3.5cm}@{}}
\toprule
 & \texttt{mean} & \texttt{covar} & \texttt{orig}\\
\midrule
Interaction written into
  & $\mu$ & $\Sigma$ & $\Sigma$\\
Interaction channel
  & BR mean shift & $b$ vector & $b$ vector\\
\addlinespace
$A_c$ within-factor
  & AR(1) $\rho_c^{|w_i-w_j|}$
  & AR(1) $\rho_c^{|w_i-w_j|}$
  & compound symmetry, constant $\rho_c$\\
$C_{cc'}$ diagonal
  & $c_1$ & $c_1$ & $c_1$\\
$C_{cc'}$ off-diagonal
  & $c_{\times}\rho^{|w_i-w_j|}$
  & $c_{\times}\rho^{|w_i-w_j|}$
  & constant $c_{\times}$\\
\addlinespace
$b$ on drug
  & $0$ & $c_{bm}$ & $c_{bm}$\\
$b$ off drug
  & $0$ & $c_{bm}e^{-\lambda t_{sd}}$ & $0$ or $c_{bm}$\\
$b$ off-drug profile
  & n/a & graded & step\\
$b$ gate
  & n/a & \texttt{onDrug}, then $t_{sd}$ & $\mu_{BR} \neq 0$\\
Gate position carryover-dependent
  & n/a & no & \textbf{yes}\\
\addlinespace
$\mu$ shift on drug
  & $c_{bm} z_i \sigma_{BR}$ & none & none\\
$\mu$ shift off drug
  & $0$ (step) & none & none\\
\addlinespace
Carryover on BR mean
  & \multicolumn{3}{l}{identical line in all three}\\
$\mathrm{BL}$ row and column
  & \multicolumn{3}{l}{identity in all three}\\
$\mathrm{BM}$--$TV$, $\mathrm{BM}$--$PB$
  & \multicolumn{3}{l}{zero in all three}\\
\bottomrule
\end{tabular}
\end{center}

Three features of this comparison are worth drawing out.

First, \texttt{mean} and \texttt{covar} share their entire covariance
matrix apart from the vector $b$. Setting $b$ to zero and adding the
mean shift converts one into the other, with every within-factor and
cross-factor block untouched. This is the decomposition of Appendix A
made concrete: the two architectures are not different models but the
two places a single joint distribution can carry an interaction.

Second, \texttt{orig} differs from \texttt{covar} on three rows
simultaneously, in $A_c$, in the off-diagonal of $C_{cc'}$, and in
$b$. No single row can therefore be held responsible for any
difference in results obtained under it, and a comparison against
\texttt{orig} bounds the joint effect of the three differences rather
than decomposing them. Separating them would require running the
correlation structure and the correlation gate as crossed factors.

Third, only the $b$ row differs in mechanism rather than in decay
profile. The $A_c$ and $C_{cc'}$ rows describe the same quantity
decaying at different rates, or not decaying at all. The $b$ row
describes different conditions being tested to decide whether the
interaction is present. That distinction is what makes the gate row
consequential: because \texttt{orig} tests the response mean rather
than the drug indicator, and carryover keeps the response mean
nonzero, its gate opens at off-drug occasions precisely when carryover
is present, and its interaction channel loses the on-drug versus
off-drug contrast on which identification depends.

The final three rows record what all three processes hold in common.
The baseline term is uncorrelated with everything, the biomarker
couples to no factor other than the response, and the carryover
adjustment to the response mean is implemented by the same line of
code in every case.

\section{Summary of differences}

\begin{center}
\begin{tabular}{lll}
\toprule
Component & \texttt{orig} & \texttt{covar}\\
\midrule
$A_c$ interior & $\rho_c$ & $\rho_c^{|w_i - w_j|}$\\
$C_{cc'}$, $i = j$ & $c_1$ & $c_1$\\
$C_{cc'}$, $i \neq j$ & $c_{\times}$ & $c_{\times}\rho^{|w_i-w_j|}$\\
$b_p$, on drug & $c_{bm}$ & $c_{bm}$\\
$b_p$, off drug & $0$ or $c_{bm}$ & $c_{bm}e^{-\lambda t_{sd,p}}$\\
$\bl$ row/column & identity & identity\\
$\bm$--$TV$, $\bm$--$PB$ & $0$ & $0$\\
\bottomrule
\end{tabular}
\end{center}

The compound-symmetry form constrains $\rho_c$ more tightly than the
AR(1) form as $n$ grows, which is the mechanism behind the narrower
positive-definite range of $c_{bm}$ under \texttt{orig}.

\end{document}
